\section{Related Work}

\subsection{Medical Retrieval-Augmented Question Answering}

Recent benchmarking studies have demonstrated that retrieval quality dominates downstream medical QA performance---often more than the choice of generator model~\citep{xiong2024benchmarking}. The MedRAG/MIRAGE benchmark systematically evaluated corpora, retrievers, and backbone models, finding that even strong generators degrade substantially when fed weak evidence. Iterative medical RAG addresses this by generating follow-up retrieval queries~\citep{xiong2024imedrag}, but the core evidence selection step remains pairwise---each passage is scored independently against the query. What these systems share is a focus on \emph{using} retrieved evidence for generation; what they lack is a mechanism for reranking evidence using the higher-order biomedical relations that connect passages to each other and to structured knowledge.

\subsection{Biomedical QA Datasets and Semantic Retrieval}

BioASQ~\citep{tsatsaronis2015bioasq} and its manually curated QA extension BioASQ-QA~\citep{krithara2023bioasqqa} provide the primary evaluation framework for biomedical evidence retrieval, pairing questions with gold passage-level annotations. PubMedQA~\citep{jin2019pubmedqa} and BEIR's NFCorpus~\citep{thakur2021beir} offer complementary testbeds with different evidence structures. For biomedical semantic retrieval, MedCPT~\citep{jin2023medcpt} trains a dual-encoder retriever and a cross-encoder reranker on large-scale PubMed search logs, establishing strong neural baselines. However, these semantic models operate on query-passage pairs in isolation. They exploit no structured biomedical knowledge---MeSH hierarchies, biomedical entity graphs, or curated relation databases---even though such resources are freely available for the biomedical domain. This gap motivates our work: can structured biomedical knowledge, when added to retrieval features, further improve evidence reranking?

\subsection{Graph and Hypergraph Approaches to RAG}

A growing line of work applies graph structure to retrieved information. GraphRAG~\citep{edge2024graphrag} organizes passages into entity-relationship graphs and performs local-to-global search for summarization. Medical Graph RAG~\citep{wu2024medgraphrag} grounds medical responses by linking user documents to credible sources and controlled vocabularies. Hypergraph extensions argue that n-ary relations---which connect more than two nodes at once---are more natural for representing complex information structures. HyperGraphRAG~\citep{luo2025hypergraphrag} uses hyperedges for passage clustering, and cross-granularity HGRAG~\citep{wang2026hgrag} models entities and passages at different retrieval granularities. Medical HyperRAG~\citep{xu2026medicalhyperrag} constructs a ClinBridge hypergraph over clinical records, cases, and guidelines for patient-centered QA.

A common thread runs through these approaches: they apply graph structure \emph{without domain-grounded concept annotations}. Edges are typically built from lexical overlap, embedding similarity, or co-occurrence statistics. Our work tests the hypothesis that this is insufficient---that raw topology, without medical knowledge to constrain it, can introduce ranking noise rather than signal. By systematically comparing the same hypergraph architecture with and without MeSH, entity, and PrimeKG annotations, we provide direct evidence that domain knowledge constraints are necessary for structural features to be useful in biomedical evidence reranking.

\subsection{Controlled Biomedical Knowledge}

MeSH~\citep{nlm2026mesh} is a hierarchically organized controlled vocabulary maintained by the National Library of Medicine, providing standardized annotations for indexing biomedical literature. PrimeKG~\citep{chandak2023primekg} integrates multimodal precision-medicine relations across diseases, proteins, drugs, and phenotypes. Both resources are publicly available and underpin biomedical information retrieval systems, yet neither has been systematically integrated into learning-to-rank pipelines for evidence retrieval. Our framework leverages MeSH hierarchy features as the primary structured knowledge signal---accounting for 75.9\% of evidence rescues---and treats PrimeKG as an auxiliary source whose exact-name coverage limits its independent contribution.
